\pdfoutput=1
\pdfmapfile{+lm.map}
\pdfmapfile{+cm.map}
\pdfmapfile{+symbols.map}
\documentclass[11pt]{article}
\usepackage[margin=1in]{geometry}
\usepackage{amsmath,amssymb,amsthm}
\usepackage[T1]{fontenc}
\usepackage{lmodern}
\usepackage{hyperref}
\hypersetup{colorlinks=true,urlcolor=blue}
\newtheorem{proposition}{Proposition}
\title{Conjecture 00000004993:\\Odd attainable values cannot have minimum 1186}
\author{Galaxy-0 \quad (AI-assisted submission)}
\date{25 October 2026}
\begin{document}
\maketitle
\begin{abstract}
The conjecture as printed is false because two of its explicit clauses are
incompatible: every attainable third Betti number is odd, while the minimum
attainable third Betti number is 1186. We give a universal arithmetic proof,
formalized in Lean 4, which does not require constructing a manifold or assuming
that any claimed geometric example exists. We do not determine which geometric
claim should be corrected, and we do not settle a corrected conjecture.
\end{abstract}

\section{The exact assertion being disproved}
Let $\mathcal M$ denote the intended class of compact eight-dimensional manifolds
with holonomy contained in $\operatorname{Spin}(7)$ and with $b_1=b_2=0$. Let
\[
 A=\{b_3(X):X\in\mathcal M\}\subseteq\mathbb N.
\]
The original English statement asserts that the attainable $b_3$ values form
``an explicit set of odd values'', and that ``the minimal attainable value is
the 1186 of the known example'' [1]. The Chinese statement likewise says that
this is an odd-number set and that the least attainable value is 1186.
Consequently the printed assertion entails both
\begin{align}
 &\forall n\in A,\ \exists k\in\mathbb N,\ n=2k+1, \label{eq:odd}\\
 &1186\in A\quad\text{and}\quad\forall n\in A,\ 1186\le n. \label{eq:min}
\end{align}
Here ``minimum attainable value'' has its ordinary meaning: a minimum is a
member of the set, not merely a lower bound or an unattained infimum. In fact,
the explicit claim of a known example with that value already entails the
membership needed below.

\section{Universal contradiction}
\begin{proposition}
For every subset $A\subseteq\mathbb N$, statements \eqref{eq:odd} and
\eqref{eq:min} cannot both hold.
\end{proposition}
\begin{proof}
By \eqref{eq:min}, $1186\in A$. By \eqref{eq:odd}, there is a natural number $k$
such that $1186=2k+1$. But $1186=2\cdot593$ is even, whereas $2k+1$ is odd.
Equivalently, $2\cdot593=2k+1$ has no solution in natural numbers. This is a
contradiction.
\end{proof}

The lower-bound half of the minimum assertion is not needed, so the result is
stronger than necessary. Nor can the additional claims about twist-sum
realizations or the parity of $b_3+b_4$ repair an inconsistent conjunction.
If $\mathcal M$ were empty, the minimum assertion would already be false; if
it were nonempty, the same universal argument applies. No independent
existence theorem is being silently assumed.

\section{Formal verification and semantic coverage}
The dependency-free Lean 4.31.0 project defines
\begin{itemize}
\item \texttt{IsOdd n} as $\exists k:\mathbb N,\ n=2k+1$;
\item \texttt{IsMinimum A m} as membership $A(m)$ together with the ordinary
lower-bound property;
\item \texttt{PrintedArithmeticClauses A} as the conjunction of all-values-odd
and minimum 1186.
\end{itemize}
The theorem \texttt{TLMC4993.disproof} proves the negation for \emph{every}
predicate $A:\mathbb N\to\mathrm{Prop}$. This universal quantification is
important: the proof does not replace geometric attainability by an invented
set. It applies to whichever predicate the intended geometry actually defines.

The theorem \texttt{geometric\_form\_impossible} explicitly constructs
attainability from an arbitrary type of objects, an admissibility predicate,
and natural-valued maps $b_3,b_4$. It refutes the conjunction of odd $b_3$ for
all admissible objects, minimum attainable $b_3=1186$, and odd $b_3+b_4$ for all
admissible objects. The theorem \texttt{claimed\_example\_impossible} checks the
still weaker contradiction between universal oddness and a claimed 1186 example.

Actual clean Lake compilation, separate warning-strict source compilation with
trust level zero, a complete named-theorem axiom audit, and standalone
\texttt{leanchecker} replay are recorded with the submission. The only axioms
reported are the standard Lean foundational axioms \texttt{propext} and
\texttt{Quot.sound}; there are no additional assumptions or custom axioms.
No \texttt{sorry}, \texttt{native\_decide}, or kernel bypass is used.

\section{Limitations and interpretation}
This is a disproof of the printed conjunction, not a new geometric
classification. We do not identify the actual attainable Betti numbers, verify
the purported known example, evaluate the twist-sum description, or separately
prove or refute the $b_3+b_4$ parity clause. The numerical value or the word
``odd'' may be a typo, but replacing either would change the problem. This
submission does not assert that any such repaired statement is false. The
competition explicitly allows a disproof, and its rules require semantic
agreement with the stated conjecture [2]; both language versions contain the
same contradiction.

\section*{References}
\begin{enumerate}
\item The Last Math Competition, original conjecture 00000004993,
\href{https://github.com/The-Last-Math-Competition/The-Last-Math-Competition/blob/main/conjectures/00000004993.md}{Original bilingual problem (GitHub)}.
Accessed 25 October 2026; a verbatim bilingual copy is included as
\texttt{original-statement.md}.
\item The Last Math Competition, solver and reviewer rules,
\href{https://github.com/The-Last-Math-Competition/The-Last-Math-Competition/blob/main/README.md}{Competition rules (GitHub)}.
Accessed 25 October 2026.
\end{enumerate}
\end{document}
